\subsection{Example E5: \texorpdfstring{$B^*$}{B*} ladder worked table}\label{ex:e5}
\paragraph{Depends on.} \Cref{def:b-extraction,def:bstar-histogram,choice:high-rung-pivot,rem:jensen,def:attenuation-operator,hook:h7}.
\paragraph{Local readout.} The same motif-conditioned closure stack under three declared return ladders.
\paragraph{Native output.}
{\footnotesize
\begin{center}
\begin{tabular}{@{}lllll@{}}
\toprule
Row & $b(\kappa)$ samples & $H_B$ support & $B^*$ & Note \\
\midrule
E5.A & $(2,2,2,4)$ & $H_B(2)=3,\ H_B(4)=1$ & $2$ & minimal two-cycle ladder \\
E5.B & $(5,5,5,7,7)$ & $H_B(5)=3,\ H_B(7)=2$ & $5$ & four-cycle enclosure ladder \\
E5.C & $(137,137,139)$ & $H_B(137)=2,\ H_B(139)=1$ & $137$ & high-rung attractor witness \\
\bottomrule
\end{tabular}
\end{center}}
\paragraph{Jensen discipline.} Histogram weighting and attenuation must be formed on the cited sample set before any averaging. If $T_{\mathrm{eval}}=n_1+n_2+n_3$ and $\bar m=(n_1+2n_2+3n_3)/T_{\mathrm{eval}}$, then the exact native loss-average is
\[
\frac{n_1+n_2\log_3\!\left(\frac{3}{2}\right)}{T_{\mathrm{eval}}},
\]
whereas the post-averaged surrogate is $\log_3(3/\bar m)$. Their difference
\[
J_{\mathrm{gap}}(\omega):=
\frac{n_1+n_2\log_3\!\left(\frac{3}{2}\right)}{T_{\mathrm{eval}}}-\log_3\!\left(\frac{3}{\bar m}\right)\ge 0
\]
is the declared Jensen-gap witness on the same window.
